\ExamPart{Trắc nghiệm đúng sai}{2}{Thí sinh trả lời từ câu 1 đến câu 2. Trong mỗi ý a), b), c), d) ở mỗi câu học sinh chọn đúng hoặc sai.}

\NQuestion Từ các chữ số $0,1,2,3,4$, xét các khẳng định sau:
\begin{SubQuestions}
  \TFItem{false}{Có $60$ số tự nhiên có $3$ chữ số đôi một khác nhau.}
  \TFItem{true}{Có $18$ số tự nhiên lẻ có $3$ chữ số đôi một khác nhau.}
  \TFItem{true}{Có $12$ số tự nhiên có $3$ chữ số đôi một khác nhau và chia hết cho $5$.}
  \TFItem{false}{Có $28$ số tự nhiên chẵn có $3$ chữ số đôi một khác nhau.}
\end{SubQuestions}
\SolutionBlock{\begin{itemize}[leftmargin=1.5em]
\item Tổng số số tự nhiên có $3$ chữ số đôi một khác nhau là $4\cdot4\cdot3=48$, nên ý a) sai.
\item Số lẻ: chọn hàng đơn vị là $1$ hoặc $3$ có $2$ cách, chọn hàng trăm có $3$ cách, chọn hàng chục có $3$ cách, được $2\cdot3\cdot3=18$, nên ý b) đúng.
\item Chia hết cho $5$ thì tận cùng là $0$. Hàng trăm có $4$ cách, hàng chục có $3$ cách, nên được $12$ số. Ý c) đúng.
\item Số chẵn gồm hai trường hợp: tận cùng là $0$ có $4\cdot3=12$ số; tận cùng là $2$ hoặc $4$ có $2\cdot3\cdot3=18$ số. Tổng cộng $30$ số, nên ý d) sai.
\end{itemize}}

\NQuestion Trong mặt phẳng tọa độ $Oxy$, cho $A(1;2),\ B(5;4),\ C(-1;-2)$. Xét các khẳng định sau:
\begin{SubQuestions}
  \TFItem{true}{$\overrightarrow{AB}=(4;2)$.}
  \TFItem{true}{Trung điểm của đoạn thẳng $BC$ là $I(2;1)$.}
  \TFItem{true}{Đường thẳng $AB$ có phương trình $x-2y+3=0$.}
  \TFItem{false}{Khoảng cách từ điểm $C$ đến đường thẳng $AB$ bằng $\dfrac{4}{\sqrt5}$.}
\end{SubQuestions}
\SolutionBlock{\begin{itemize}[leftmargin=1.5em]
\item $\overrightarrow{AB}=(5-1;4-2)=(4;2)$, nên ý a) đúng.
\item Trung điểm của $BC$ là $
I\left(\dfrac{5+(-1)}{2};\dfrac{4+(-2)}{2}\right)=(2;1),$ nên ý b) đúng.
\item Hệ số góc của $AB$ là $\dfrac{4-2}{5-1}=\dfrac12$. Phương trình qua $A(1;2)$ là $
y-2=\dfrac12(x-1)\Leftrightarrow x-2y+3=0,$ nên ý c) đúng.
\item Khoảng cách từ $C(-1;-2)$ đến $AB:x-2y+3=0$ là $
d(C,AB)=\dfrac{|(-1)-2(-2)+3|}{\sqrt{1^2+(-2)^2}}=\dfrac{6}{\sqrt5},$ không phải $\dfrac{4}{\sqrt5}$, nên ý d) sai.
\end{itemize}}
